\section{Content Moderation Schemes}
\label{sec:content_mod}

%In this section, we will describe the content moderation schemes realizing \contentModFuncExact and \contentModFuncFuzzy leveraging the blocklist constructions described in \secref{sec:acb}. 


\subsection{Exact Matching: Realizing \contentModFuncExact}
%
We describe our content moderation scheme with exact-matching capabilities in \algoref{protocol:em}; we assume that the client wants to upload content drawn from the universe \universe, and the content hash function $\contentHash: \universe \rightarrow \bin^{\genericVecLen}$ maps any such content to an \genericVecLen-bit string. Two items $\genericUniverseElmt, \genericUniverseElmtAlt$ are similar iff. $\contentHash(\genericUniverseElmt) = \contentHash(\genericUniverseElmtAlt)$.  


\myparagraph{Key idea: fast path with two ACBs}
%
We can directly leverage either a Bloom filter based blocklist or a signature based blocklist to realize \contentModFuncExact. In the former, the cost of matching against the blocklist is $\bigO(1)$ but the construction has false positives. In the latter, there are no false positives (or false negatives) but the costs for matching is linear (in the blocklist size). This motivates the questions: \emph{can we design a scheme that has sublinear matching costs without incurring false positives?} Our solution is to combine the two ACBs to realize \contentModFuncExact without false positives and sublinear matching costs (\algoref{protocol:em}). 


\noindent\textbf{Creating the Blocklist:}
\enumListStart

\item The auditors partition the blocklisted items into \numBlocklistParts disjoint partitions $\blocklistPartition{1}, \dots, \blocklistPartition{\numBlocklistParts}$. 

\item The auditors invoke $\sigACB.\publishList(\blocklistPartition{1}, \dots, \blocklistPartition{\numBlocklistParts})$ and publish a separate signature-based ACB for each of the blocklist partitions $\sigBlklistToken{1}, \dots, \sigBlklistToken{\numBlocklistParts}$ 

\item The auditors invoke $\bfACB.\publishList(\blocklist)$ and publish a blinded blocklist \blindedBlocklist

\item The client and server invoke $\bfACB.\preprocess(\blindedBlocklist)$ and the client learns \bfBlklistToken

\enumListEnd

\noindent\textbf{Upload:}
When uploading some content \clientData
\enumListStart
\item (\emph{fast path}:) The client derives the content hash $\genericCInputElmt \assign \embed(\clientData)$ and invokes $\bfACB.\searchBlocklist(\genericCInputElmt, \bfBlklistToken)$. If the search returns 0, then the client and server can abort because this means $\genericCInputElmt \notin \blocklist$. The server will store the encrypted content. 

\item (\emph{fine-grained check}:) Otherwise, let $\bloomFilter[\bloomFilterIdx]$ be the Bloom filter such that $\prp{\prpKeyServer}(\prf{\prfKeyClient}(\genericCInputElmt)) \in \bloomFilter[\bloomFilterIdx]$; recall that the server learns $\bloomFilter[\bloomFilterIdx]$ in $\bfACB.\searchBlocklist$. The server identifies the blocklist partitions containing the items included in $\bloomFilter[\bloomFilterIdx]$. More specifically, the server calculates the set of indices $$\partitionIDSet \assign \{ \partitionIdx: \genericBlocklistElmt \in \bloomFilter[\bloomFilterIdx] ~\land ~\genericBlocklistElmt \in \blocklistPartition{\partitionIdx}\}$$, and sends it to the client 


\item Upon receiving $\partitionIDSet$, the client invokes $\sigACB.\searchBlocklist(\genericCInputElmt, \sigBlklistToken{\partitionIdx})$ for all $\partitionIdx \in \partitionIDSet$. If any invocation returns $\genericCInputElmt$ to the server, then $\genericCInputElmt \in \blocklist$; the server rejects the encrypted content and aborts. Otherwise, $\genericCInputElmt \notin \blocklist$ and the server accepts the content and informs the client.  
\enumListEnd

\myparagraph{Costs}
%

\itemListStart
\item The preprocessing for \bfBlklistToken requires a one-time $\bigO(\blocklistSize)$ compute and communication cost for the client and the server. The client also requires $\bigO(\blocklistSize +
 \numBlocklistParts)$ storage. 
 
\item During an upload, $\bigO(1)$ compute cost for the client and server (with high probability) when the uploaded content does not match an item in the blocklist, i.e., fast path verification. Alternatively, $\bigO(\setSize{\partitionIDSet}) \in \bigO(\log \blocklistSize)$ compute cost for the client and $\bigO(\blocklistSize)$ server cost when the client's uploaded file matches an item on the blocklist
\item $\bigO((\blocklistSize/\numBlocklistParts)+ \log \blocklistSize)$ communication cost for sending an encrypted Bloom filter, and $\bigO(\log \blocklistSize)$ randomized signatures if a fine-grained check is deemed necessary. If $\numBlocklistParts \in \bigO(\blocklistSize)$, then the cost of an upload is $\bigO(\log \blocklistSize)$. 
\itemListEnd



\begin{algorithm}[t]
\caption{\protocolEM: content moderation with exact matching}
\label{protocol:em}
\footnotesize
\begin{algorithmic}[1]

\Statex\LComment{Partition the blocklist and publish signature-based ACB}
\State Auditors distribute \blocklist into \numBlocklistParts disjoint buckets $\blocklistPartition{1}, \dots, \blocklistPartition{\numBlocklistParts}$ 

\For{each $\partitionIdx \in [1, \numBlocklistParts]$} 
\State Auditors invoke $\sigACB.\publishList(\blocklistPartition{\partitionIdx})$ and publish the signature-based ACB $\sigBlklistToken{\partitionIdx} = \langle \auditorPubKeyOurScheme, \genericSig[\blocklistPartition{\partitionIdx}] \rangle$
\State Auditors sign and publish $\{ \signGenSig[\auditorPrivKeyGenericScheme[\auditorIdx]](\hashFn(\sigBlklistToken{\partitionIdx})) \}_{\auditorIdx \in [1, \numAuditors]}$
\EndFor

\Statex\LComment{Publish blinded blocklist for Bloom filter based ACB}
\State Auditors invoke $\bfACB.\publishList(\blocklist)$ and publish a blinded blocklist $\blindedBlocklist$ and the signatures $\{ \signGenSig[\auditorPrivKeyGenericScheme[\auditorIdx]](\hashFn(\blindedBlocklist) ) \}_{\auditorIdx \in [1, \numAuditors]}$

\Statex\LComment{Client and server preprocess blinded blocklist}

\State \client downloads $\blindedBlocklist$,  $\{ \signGenSig[\auditorPrivKeyGenericScheme[\auditorIdx]](\hashFn(\blindedBlocklist)) \}_{\auditorIdx \in [1, \numAuditors]}$, $\{\sigBlklistToken{\partitionIdx}\}_{\partitionIdx \in [1, \numBlocklistParts]}$, $\{ \{ \signGenSig[\auditorPrivKeyGenericScheme[\auditorIdx]](\hashFn(\sigBlklistToken{\partitionIdx})) \}_{\auditorIdx \in [1, \numAuditors]}\}_{\partitionIdx \in [1, \numBlocklistParts]}$, and verifies the signatures. If the signatures do not verify, then \client aborts. 

\State \client and \server invoke $\bfACB.\preprocess(\blindedBlocklist)$. \client stores $\bfBlklistToken \assign \{\encrypt{\genericSymmKey}(\bloomFilter[1]), \dots, \encrypt{\genericSymmKey}(\bloomFilter[\numBins])\}$ 


\Statex\LComment{Matching against blocklist during uploads}

\State \client sends $\langle \genericFileID, \encrypt{\fileKey}(\clientData) \rangle$ to \server

\State \client computes $\genericCInputElmt \assign \contentHash(\clientData)$ and $\bfACB.\searchBlocklist(\genericCInputElmt, \bfBlklistToken)$. 

\If{$\bfACB.\searchBlocklist(\genericCInputElmt, \bfBlklistToken) = 0$} \Comment{Fast path}
\State \server outputs $\langle \genericFileID, \detected = \fals, \emptyset \rangle$ and aborts. 
\Else  \Comment{Fine-grained check}
\State \server computes $\partitionIDSet \assign \{ \partitionIdx: \genericBlocklistElmt \in \bloomFilter[\bloomFilterIdx] ~\land ~\genericBlocklistElmt \in \partition{\partitionIdx}\}$ and sends to \client.  
\For{each $\partitionIdx \in \partitionIDSet$}
\State \client invokes $\sigACB.\searchBlocklist\left(\genericCInputElmt,\sigBlklistToken{\partitionIdx}\right)$
\If{$\sigACB.\searchBlocklist\left(\genericCInputElmt,\sigBlklistToken{\partitionIdx}\right) = \genericCInputElmt$}
\State \server outputs $\langle \genericFileID, \detected = \tru, \genericCInputElmt \rangle$ and aborts
\Else
\State \server outputs $\langle \genericFileID, \detected = \fals, \emptyset \rangle$. \client adds $\genericCInputElmt$ to $\prevHashList$. 
\EndIf
\EndFor
\EndIf


\end{algorithmic}
\end{algorithm}

\begin{theorem}
Assuming there are protocols securely realizing \distSigningFunc and \authOPRFFunc, \protocolEM securely realizes \contentModFuncExact 
\end{theorem}

The proof is available in the full version~\cite{code}


\subsection{Fuzzy Matching: Realizing \contentModFuncFuzzy}
%
We will describe a content moderation scheme realizing \contentModFuncFuzzy. We will assume that the set of blocklisted items are embedded into a metric space with distance \distance; the embedding produces a \genericVecLen-dimensional vector, and so we will refer to 
the blocklisted items as vectors henceforth. 
Our protocol will detect matches of the form $\tupleDefn{\blocklistVec, \clientVec}$ such that 
$\distanceNotation[\distance][\genericBlocklistElmt][\genericCInputElmt] \le \threshold$. In return, the protocol will have a small number of false negatives -- cases where items are blocked but are not detected. Our protocol operates over bit-vectors in the Hamming space;  our choice of Hamming distance is informed by the fact that several image embeddings use Hamming distance to identify similar images~\cite{singh2025perceptual,neuralHashAttack,yang2025certphash} and existing work on encrypted content moderation also use Hamming distance as the metric~\cite{kulshrestha2021identifying}.


%\emph{All existing content moderation schemes that perform fuzzy matching suffer from false negatives~\cite{kulshrestha2021identifying}}. The protocol does not have false positives. 




\myparagraph{Key Idea: Using projections for a threshold match}
%
Comparing an input against the blocklist is challenging when we need to implement a fuzzy match. The problem is that to accurately detect a match, we must compute $\distanceNotation[\distance][\genericBlocklistElmt][\genericCInputElmt]$ for $\genericBlocklistElmt \in \blocklist$. One solution is to use a fuzzy membership test (PMT) or a fuzzy private set intersection(PSI) protocol~\cite{}, but they have high computational costs, and cannot scale to the latency requirements of a file server. 

We take a different approach by leveraging the fact that our functionality has strictly weaker requirements than a fuzzy PSI/PMT protocol because: i) it tolerates false positives and false negatives, and ii) it does not output the matching items to the server/client. Thus, we can design a protocol with lower costs. The key idea is to use projections of the blocklist vectors as the matching criteria: two vectors are considered close if they have identical projections under a certain number of projection functions. Obviously, leaking the projections themselves will allow the server to learn information about the encrypted content. Therefore, we only reveal the \emph{number of matching projections} when they are above a certain threshold, but not the projections themselves. Thus, even if there is a false positive, the server does not learn any more information about the content.

We do not provide a mechanism to disambiguate false positives from true positives. A fine-grained check without false positives must essentially solve the fuzzy, ``distribution-free''~\cite{} private membership problem against malicious receivers with sublinear communication costs. To the best of our knowledge, there is no such protocol that scales to the throughput requirements of a file server. Hence, we adopt a best-effort approach focusing on efficiency rather than accuracy: we will select projection functions carefully such that the chance of false positives are low. Given the false positive rate, the server can fix a threshold and decide to report/suspend an account if there are more than a threshold number of potential matches, which is likely to include a true positive. To realize this idea, we will first define the projection functions, and then show how to combine it with our cuckoo hash based ACB to realize \contentModFuncFuzzy. 

\iffalse 
, and if required, follow it by performing  a significantly more expensive but perfectly accurate test over the actual vectors. Specifically, if the client's input matches a certain threshold number of projections against the blocklisted items, then the server performs the fine-grained check., and the likely candidates in the blocklist that should be checked against the input. We note that the \contentModFunc functionality allows this leakage, since the server does not learn the content key.  Nonetheless, , and iii) the candidate set is reasonably large and can be set as a tunable parameter. A caveat here is that the fine-grained check must be performed over the actual embedding vectors and not just the projections, so the blocklist authentication token and the content verification token must be modified to account for entire vectors. , and then describe how to modify the tokens. 
\fi 

\subsubsection{Locality-Preserving Fingerprints}


\begin{definition}
	For a metric space \metricSpace, a \projectionPrimitive of an element $\genericMetricSpaceElmt$ in the set $\blocklist \subseteq \metricSpace$ 
	comprises its projections under uniformly sampled functions $\projFn[1], 
	\dots, \projFn[\numProjections]: \metricSpace \rightarrow \projectedSpace$, $\fingerprint(\genericMetricSpaceElmt) \assign 
	\{\hashFn(\projFn[1](\genericMetricSpaceElmt), \projFn[1])$ $\dots, \hashFn(\projFn[\numProjections](\genericMetricSpaceElmt), \projFn[\numProjections])\}$, where $\hashFn: \bin^{\ast} \rightarrow \setDefn{0,1}^{\secParam}$ is a collision-resistant hash function, such that given any 
$\genericMetricSpaceElmtAlt \in \metricSpace$,
		
		\begin{small}
		\begin{align}
			& \cprob{}{\not\exists \projIdx \in \nats[1][\numProjections], 
				\projFn[\projIdx](\genericMetricSpaceElmt) = \projFn[\projIdx](\genericMetricSpaceElmtAlt)}{\distanceNotation[\distance][\genericMetricSpaceElmt][\genericMetricSpaceElmtAlt] \le 
			\threshold} \le \fnr^{\numProjections} \label{eq:no_false_neg} \\
			& \cprob{}{\not\exists \projIdx \in \nats[1][\numProjections], 
				\projFn[\projIdx](\genericMetricSpaceElmt) \neq \projFn[\projIdx](\genericMetricSpaceElmtAlt)}{\distanceNotation[\distance][\genericMetricSpaceElmt][\genericMetricSpaceElmtAlt] > 
			\threshold} \le \fpr^{\numProjections} \label{eq:no_false_pos}
		\end{align}
		\end{small}
where,  \fpr and \fnr are the false positive rate and false negative rate under each projection function. 
		%Here, \tpr and \tnr are tunable true positive rate and true negative rate. 
		
%		\item For every pair of element, $\genericMetricSpaceElmt, \genericMetricSpaceElmtAlt \in \blocklist$ 
%		$\setSize{\fingerprint{\genericMetricSpaceElmt} ~\triangle ~\fingerprint{\genericMetricSpaceElmtAlt}} = 0$ with probability  
%		negligible in some parameter \secParam, where $\triangle$ is the symmetric difference. That is, there is at least one projection such that 
%		$\projFn[\projIdx](\genericMetricSpaceElmt) \neq \projFn[\projIdx](\genericMetricSpaceElmtAlt)$ with overwhelming probability. 

%		$\genericBlocklistElmt \in \blocklist$, there is at least one function $\projFn[\projIdx]$ where 
%		$\projFn[\projIdx](\genericBlocklistElmt) \notin 
%		\setDefn{\projFn[\projIdx](\genericBlocklistElmt[1]'), 
%			\dots, \projFn[\projIdx](\genericBlocklistElmt[\blocklistSize]')}$ for all 
%		$\genericBlocklistElmt[1]', 
%		\dots, \genericBlocklistElmt[\blocklistSize]' \in \blocklist\setminus\genericBlocklistElmt$ (with 
%		overwhelming probability).
		
%	\end{enumerate}

\end{definition}

Intuitively, \eqnref{eq:no_false_neg} ensures that two items close in the metric space agree on at least one of the projections 
with high probability. \eqnref{eq:no_false_pos} implies that two items far apart in the metric space do not agree on \emph{all} their projections 
with high probability. Note that \fpr and \fnr may not be negligible. Also, since the projection functions are selected randomly, the overall false positives and false negatives over the projections are binomially distributed. Specifically, 


\def\scalingFn{\ensuremath{\alpha}\xspace}

\begin{corollary}
\label{corr:fingerprint_intersection}
Let $\threshold \in \reals^{++}$ be a distance threshold, and $\scalingFn: \reals^{++} \rightarrow  [1, \numProjections]$ is a scaling function. Also, let $\detectionThreshold \assign \scalingFn(\threshold)$. Given two elements $\genericMetricSpaceElmt, \genericMetricSpaceElmtAlt \in \metricSpace$ and a detection threshold, 

\begin{gather*}
\prob{\badFN} = 
\cprob{}{\begin{array}{c} |\fingerprint{}(\genericMetricSpaceElmt) \cap \fingerprint{}(\genericMetricSpaceElmtAlt)|  \\ < \numProjections - \detectionThreshold + 1 \end{array}}{\distanceNotation[\distance][\genericMetricSpaceElmt][\genericMetricSpaceElmtAlt] \le \threshold} 
\\  \le 1 - \sum\limits_{i = 0}^{\numProjections - \detectionThreshold} \binom{\numProjections}{i} \fnr^{i} (1 - \fnr)^{\detectionThreshold - i} \\
\prob{\badFP} = \cprob{}{\begin{array}{c} |\fingerprint{}(\genericMetricSpaceElmt) \cap \fingerprint{}(\genericMetricSpaceElmtAlt)| \\ \ge  \detectionThreshold \end{array}}{\distanceNotation[\distance][\genericMetricSpaceElmt][\genericMetricSpaceElmtAlt] > \threshold} 
\\ \le \sum\limits_{i = \detectionThreshold}^{\blocklistSize \numProjections} \binom{\blocklistSize\numProjections}{i} \fpr^{i} (1 - \fpr)^{\blocklistSize \numProjections - i} \\
\end{gather*}

\end{corollary}

The proof is available in the full version~\cite{code}. 


\iffalse 
\begin{corollary}
 \label{eq:false_pos_bound} 
Let $\detectionThreshold \in [1, \numProjections]$. Given a dataset $\blocklist$ of $\blocklistSize$ elements, all having unique projections under $\numProjections$ random projection functions $\projFn[1], \dots, \projFn[\numProjections]$, and a  test point $\genericMetricSpaceElmt$. The probability that for a subset of projection functions $I$, where $\setSize{I} \ge \detectionThreshold$, the probability that $\projFn[\projIdx](\genericMetricSpaceElmt) = \projFn[\projIdx](\genericMetricSpaceElmt')$ for some $\genericMetricSpaceElmt' \in \blocklist$ and $\distanceNotation[\distance][\genericMetricSpaceElmt][\genericMetricSpaceElmtAlt] > 
			\threshold$ is bounded by 


\[			 \prob{\badFP}
 \le \sum\limits_{i = \detectionThreshold}^{\blocklistSize \numProjections} \binom{\blocklistSize\numProjections}{i} \fpr^{i} (1 - \fpr)^{\blocklistSize \numProjections - i}
\]
, where  \metricSpace is a metric space with distance $\distanceNotation[\distance][\cdot]$, \fpr is the false positive rate under each random projection.
\end{corollary}


\begin{corollary}
 \label{eq:false_neg_bound}
Let $\detectionThreshold \in [1, \numProjections]$ and two elements $\genericMetricSpaceElmt, \genericMetricSpaceElmtAlt \in \metricSpace$ such that $\hammingDist{\genericMetricSpaceElmt}{\genericMetricSpaceElmt} \le \threshold$. Then the probability that they do not agree on at least $\numProjections - \detectionThreshold + 1$ projections is bounded by 

\[	
\prob{\badFN} \le 1 - \sum\limits_{i = 0}^{\numProjections - \detectionThreshold} \binom{\numProjections}{i} \fnr^{i} (1 - \fnr)^{\detectionThreshold - i}
\]
, where \fnr is the false negative rate under each random projection.

\end{corollary}


\begin{corollary}
\label{corr:fingerprint_intersection}
Let $\detectionThreshold \in [1, \numProjections]$. Given two elements $\genericMetricSpaceElmt, \genericMetricSpaceElmtAlt \in \metricSpace$ and a detection threshold, 

\begin{align*}
\cprob{}{|\fingerprint{}(\genericMetricSpaceElmt) \cap \fingerprint{}(\genericMetricSpaceElmtAlt)| < \numProjections - \detectionThreshold + 1}{\distanceNotation[\distance][\genericMetricSpaceElmt][\genericMetricSpaceElmtAlt] \le \threshold} 
\\ \le \prob{\badFN} \\
\cprob{}{|\fingerprint{}(\genericMetricSpaceElmt) \cap \fingerprint{}(\genericMetricSpaceElmtAlt)| \ge  \detectionThreshold}{\distanceNotation[\distance][\genericMetricSpaceElmt][\genericMetricSpaceElmtAlt] > \threshold} 
\\ \le \prob{\badFP} \\
\end{align*}

\end{corollary}
\fi 

\myparagraph{Locality-preserving fingerprint for the Hamming distance}
%
\citet{chongchitmate2024approximate} describe a construction for a locality-preserving fingerprint for the Hamming space using uniformly sampled functions $\projFn[1], \dots, 
\projFn[\numProjections]$ where $\projFn[\projIdx]: \setDefn{0,1}^{\genericVecLen} \rightarrow 
\setDefn{0,1}^{\secParam}$. Each function specifies a random subset of bits from the input 
vector, and injectively maps it to a $\secParam$-bit string. Consider a set
$\blocklist \subseteq \setDefn{0,1}^{\genericVecLen}$.

\enumListStart
\item Create \numProjections key-value stores $\projBuckets[1], \dots, \projBuckets[\numProjections]$ where the keys are the vectors in \blocklist.
\item For each function $\projFn[\projIdx]$ and each element $\blocklistVec \in \blocklist$
\enumListStart
\item Evaluate  $\hashFn(\projFn[\projIdx](\blocklistVec), \projIdx)$ where $\hashFn: \bin^{\secParam + \log \numProjections} \rightarrow \bin^{\secParam}$ is a collision-resistant hash function. 
\item Insert $\hashFn(\projFn[\projIdx](\blocklistVec), \projIdx)$  in the key-value store $\projBuckets[\projIdx]$ at index $\projBuckets[\projIdx][\blocklistVec]$ unless $\projFn[\projIdx](\blocklistVec)$ is already there in $\projBuckets[\projIdx]$. If so, assign a random value from the range $\bin^{\secParam}$
\enumListEnd
\enumListEnd

The set $\fingerprint{\blocklistVec} \assign \setDefn{\projBuckets[1][\blocklistVec], \dots, \projBuckets[\numProjections][\blocklistVec]}$ is a locality-preserving 
fingerprint of $\blocklistVec$ assuming negligible collisions in $\hashFn$. If all the vectors in $\blocklist$ are at least $\lshConst 
\threshold$ distance apart, then $\numProjections = \bigO\left((\genericSetSize 
\lshConst)^{\frac{1}{\lshConst - 1}} (\secParam + \log \genericSetSize)\right)$ projections ensure that 
$\fpr^\numProjections$ and $\fnr^\numProjections$ are both negligible in \secParam; here $\setSize{\blocklist} = \blocklistSize$. Observe that if $\lshConst > 2$, the number of projections are sublinear in 
$\genericSetSize$. Most embedding functions ensure that similar 
items are close in the embedded space, while dissimilar items are sufficiently far apart~\cite{kulshrestha2021identifying, mckeown2023hamming}. Alternatively, nearby items can be repeatedly 
clustered to have an efficient construction~\cite{chongchitmate2024approximate}.



\begin{algorithm}[t]
\caption{\protocolFM: content moderation with exact matching\label{protocol:fuzzy}}
\footnotesize
\begin{algorithmic}[1]

\Statex\LComment{Create a CHT based blocklist using locality-preserving fingerprints}


\For{$\blocklistVec \in \blocklist$}
\State Auditors compute the the locality-preserving fingerprints $\{\fingerprint(\blocklistVec)\}_{\blocklistVec \in \blocklist}$ using \numProjections randomly sampled functions $\projFn[1], \dots, \projFn[\numProjections]$
\EndFor
\State Auditors invoke $\chtACB.\publishList\left(\{ \fingerprint(\blocklistVec) \}_{\blocklistVec \in \blocklist}\right)$ and publish \chtBlklistToken and $\{\signGenSig[\auditorPrivKeyGenericScheme[\auditorIdx]](\hashFn(\blklistToken))\}_{\auditorIdx \in [1, \numAuditors]}$ . 

\Statex\LComment{Threshold matching against blocklist during uploads}

\State \client sends $\langle \genericFileID, \encrypt{\fileKey}(\clientData) \rangle$ to \server

\State \client computes $\clientVec \assign \embed(\clientData)$ and $\fingerprint(\clientVec)$ under the projection functions $\projFn[1], \dots, \projFn[\numProjections]$. 

\State \client and \server invoke $\chtACB.\thresholdSearchBlocklist(\fingerprint(\clientVec), \chtBlklistToken, \detectionThreshold)$. 
\If{$\chtACB.\thresholdSearchBlocklist(\fingerprint{\clientVec}, \chtBlklistToken) = 1$}
\State \server adds $\genericFileID$ to flagged content list $\flagContentList$
\State \server outputs $\langle \genericFileID, \detected=\tru\rangle$. 
%\If{$\setSize{\flagContentList} \ge \flagContentThreshold$}
%\State \server deletes all files with $\genericFileID \in \flagContentList$ 
%\Else 
%\State \server stores the encrypted file.
%\EndIf
%\State \server outputs $\langle \genericFileID, \detected=\tru\rangle$. 
\Else 
\State \server outputs $\langle \genericFileID, \detected=\fals\rangle$
\EndIf


\end{algorithmic}
\end{algorithm}





\subsubsection{Protocol design}
%
A locality-preserving fingerprint provides a construction realizing \contentModFuncFuzzy (\algoref{protocol:fuzzy}). 

\enumListStart
\item The auditors compute the locality preserving fingerprints of all the blocklist vectors. 
\item The auditors publish an encrypted cuckoo hash table based blocklist by invoking $\chtACB.\publishList$ with the locality-preserving fingerprints as input
\item During an upload, 
\enumListStart
\item The client computes $\clientVec \assign \embed(\clientData)$
\item The client and server invoke $\chtACB.\thresholdSearchBlocklist(\fingerprint(\clientVec), \chtBlklistToken, \detectionThreshold)$ to perform a threshold search which returns $\tru$ when 
$|\fingerprint(\clientVec) \cap \chtBlklistToken| \ge \detectionThreshold$. Due to \corrref{corr:fingerprint_intersection}, there is some vector $\blocklistVec \in \blocklist$ such that $\hammingDist{\clientVec}{\blocklistVec} \le \threshold$ with probability at least $1 - \prob{\badFP}$. 

\item When $\chtACB.\thresholdSearchBlocklist(\fingerprint(\clientVec), \chtBlklistToken, \detectionThreshold) = \tru$, the server adds the ID to a \emph{flagged content list} $\flagContentList$ for that specific account 

%\item If there are already a predefined threshold \revealThreshold number of content IDs in \flagContentList, then the server deletes all the corresponding content and informs the client. Otherwise, the server accepts the encrypted content and stores it. 
\enumListEnd
\enumListEnd


Once the flagged content list for an account exceeds a certain threshold \revealThreshold, the server may suspend the account and/or delete the flagged files. The threshold \revealThreshold is determined using the probability of a false positive \badFP. Since the projection functions are selected randomly, the probability of multiple false positives drops exponentially. Thus, when the number of matches exceeds the expected number of false positives by a reasonable factor, the server can infer that there is a true positive with high probability. We provide an analysis in the full version~\cite{code}. 


\myparagraph{Cost \& Security }
%
Creating the ACB has $\bigO(\blocklistSize)$ cost, and as discussed in \secref{sec:acb}, the cost of a threshold search is $\bigO(\detectionThreshold)$ for the client and $\bigO(\detectionThreshold \times \numProjections)$ for the server.   
\begin{theorem}
Assuming there is an IND-CPA secure additively homomorphic encryption scheme, \protocolFM securely realizes \contentModFuncExact with false positive rate \badFP and false negative rate \badFN. 
\end{theorem}

The proof is available in the full version~\cite{code}.